\documentclass[10pt,a4paper]{amsart}
\usepackage[margin=19mm]{geometry}
\usepackage{amsmath,amssymb,mathtools}
\usepackage[scr=rsfs,cal=euler]{mathalfa}
\usepackage{xfrac}
\usepackage[unicode,hidelinks,bookmarks=false]{hyperref}
\newcommand{\ie}{i.e.\ }
\newcommand{\cf}{cf.\ }
\newcommand{\resp}{respectively\ }
\newcommand{\Subsection}{\subsection}
\providecommand{\nobreakdash}{\nobreak\hbox{-}}
\setlength{\emergencystretch}{2em}
\multlinegap0pt
\DeclareMathOperator{\PSL}{PSL}
\usepackage{amssymb}
\usepackage{xcolor}
\newtheorem{lemma}{Lemma}
\newcommand{\FF}{\mathbb F}
\newcommand{\abs}[1]{\lvert#1\rvert}
\newcommand*{\defset}[2]{\{\,#1\,\mid\,#2\,\}}
\newcommand*{\set}[1]{\{#1\}}
\title{Transitive sets of derangements in primitive actions of $\PSL_2(q)$}
\author{Peter M\"uller}
\date{Source: arXiv:2512.19500v1 (CC BY 4.0)}
\begin{document}
\maketitle

\noindent\emph{Source and licence.} Transitive sets of derangements in primitive actions of $\PSL_2(q)$. Version arXiv:2512.19500v1 (CC BY 4.0), distributed by arXiv under the Creative Commons Attribution 4.0 International licence, http://creativecommons.org/licenses/by/4.0/, as recorded in the arXiv licence metadata for this exact version and shown on the abstract page https://arxiv.org/abs/2512.19500v1. Full licence notice: \texttt{licenses/arxiv-2512.19500-CC-BY-4.0.txt}.\par\smallskip

\noindent\textit{Editorial scope.} Counted unit: the article's lemma lemma (source0.tex lines 345--351). The article's complete original proof of it is delivered with it (source0.tex lines 352--438, one \texttt{proof} environment of 87 lines). No mathematics is written, repaired or re-derived: the statement and the proof are the article's own lines, in the article's own notation, and no retained line is substituted or re-flowed.  Retained uncounted as the source-local context the proof needs: lines 225--229.  Nothing was omitted from any retained span, so the package carries no omission record.  No retained line contains a citation, so the wrapper carries no bibliography and no bibliography entry.  No retained line contains a figure environment, an image-inclusion command or drawing code, so no image is delivered and the package declares no asset dependency.  Editorial formatting only, in addition: an AMS \texttt{amsart} preamble that restates the article's own declarations and macros used by the retained lines, namely the environments \texttt{lemma} on separate counters, together with the standard packages those lines need.  The retained environments print in this document as Lemma 1 in order of appearance, whereas the article prints the same items with the numbering of its own full document.  The complete native source of the article as distributed in the arXiv e-print for this exact version is bundled unchanged as \texttt{sources/191482/source0.tex}.
\begin{lemma}\label{l:hasse}
  Let $q$ be an odd prime power and let $f\in\FF_q[X]$ be a separable
  polynomial of degree $\le4$. Then the number of pairs
  $(x,y)\in\FF_q^2$ with $y^2=f(x)$ is at most $q+1+2\sqrt{q}$.
\end{lemma}

\begin{lemma}
  Let $q$ be a prime power with $q\ge23$ if $q$ is odd. For
  $F=\FF_{q^2}$ set $N=\defset{r\in F}{r^{q+1}=1}$. Suppose that for
  $0\ne a\in F$ the following holds: For all $r\in N$ with
  $ar+a^q/r\ne0$ there is $s\in N$ such that $ar+a^q/r = s+1/s$. Then
  $a\in N$.
\end{lemma}

\begin{proof} We first handle the easier case that $q$ is even.
As the square map is an automorphism of $F$, there is $r_0\in F$ with
$a^{q-1}=r_0^2$. Note that $a^{q-1}\in N$, hence $r_0\in N$.
  
From%
\[
  ar+a^q/r=ar+ar_0^2/r=ar_0(r/r_0+r_0/r)
\]
  we see that in order to prove the lemma, we may assume $a^q=a$.

  Set $T=\defset{r+1/r}{r\in N}$. We see that multiplication by $a$
  permutes the elements of the set $T$. Thus $ax=x$ for
  $x=\sum_{t\in T}t^2$. Let $\zeta\in F$ have order $q+1$. Then the
  elements $\zeta^i+\zeta^{-i}$, $i=0,1,\dots,q/2$ are the elements
  from $T$ without repetitions. Therefore
  \[
    x = \sum_{i=0}^{q/2}(\zeta^{2i}+\zeta^{-2i})%
    =
    \frac{\zeta^{2+q}-1}{\zeta^2-1}+\frac{\zeta^{-2-q}-1}{1/\zeta^2-1}%
    =\frac{\zeta-1}{\zeta^2-1}+\frac{1/\zeta-1}{1/\zeta^2-1}=1,
  \]
  hence $a=1$ and we are done.

  Now assume that $q$ is odd. An argument as above does not work
  because if $a^{(q-1)/2}$ is not in $N$, then
  $\defset{ar+a^q/r}{r\in N}$ is a proper subset of
  $\defset{s+1/s}{s\in N}$. (Note that $ar+a^q/r$ is fixed under
  $r\mapsto a^{q-1}/r$, so this map would not have a fixed point, in
  contrast to $s\mapsto 1/s$ which fixes $s+1/s$ for $s=\pm1$.)

  We use a finite field analog of the Cayley transform from complex
  numbers to parametrize the ``circle'' $N=\defset{r}{r\cdot r^q=1}$.
  Pick $0\ne\omega\in F$ with $\omega^q=-\omega$. Note that
  $\gamma=\omega^2$ is in $\FF_q$. Set $N'=N\setminus\set{1}$. One
  quickly checks that
  \[
    \FF_q\to N', z\mapsto\frac{z-\omega}{z+\omega}
  \]
  is bijective.

  Set $f(r)=ar+a^q/r$. We want to compute a lower bound for the size
  of $M=\defset{(r,s)\in N'\times N'}{f(r)=s+1/s}$. Note that given
  $c$, there are at most two values for $r$ with $f(r)=c$. Therefore,
  there are at least $\abs{N'}-6=q-6$ choices $r\in N'$ such that
  $f(r)\notin\set{-2, 0, 2}$. For these $r$, there is $s\in N'$ with
  $(r, s)\in M$. But then $(r,1/s)\ne(r,s)$ is also in $M$. Thus
  $\abs{M}\ge2(q-6)$.

  With $r=\frac{x-\omega}{x+\omega}$ and $s=\frac{y-\omega}{y+\omega}$
  we get that
  \[
    a\frac{X-\omega}{X+\omega} + a^q\frac{X+\omega}{X-\omega} =
    \frac{Y-\omega}{Y+\omega} + \frac{Y+\omega}{Y-\omega}
  \]
  has at least $2(q-6)$ solutions $(x,y)\in\FF_q\times\FF_q$. Write
  $a=u+v\omega$ with $u,v\in\FF_q$. Then $a^q=u-v\omega$. For later
  use we note that
  $a^{q+1}=a\cdot a^q=(u+v\omega)(u-v\omega)=u^2-v^2\gamma$.

  Clearing denominators and using $\omega^2=\gamma$, we arrive at the
  curve $P(X)Y^2 = \gamma Q(X)$, where
  \begin{align*}
    P(X) &= \left(u - 1\right) X^{2} - 2 v \gamma X + (u+1)\gamma\\
    Q(X) &= \left(u + 1\right) X^{2} - 2 v \gamma X + (u-1)\gamma.
  \end{align*}
  Note that $P, Q\in\FF_q[X]$.

  The discriminant of $P(X)$ and $Q(X)$ is
  $-4\gamma(u^2-v^2\gamma - 1)=-4\gamma(a^{q+1}-1)$ and
  $4\gamma(u^2-v^2\gamma - 1)=-4\gamma(a^{q+1}-1)$,
  respectively. Furthermore, the resultant of $P(X)$ and $Q(X)$ is
  $16\gamma^2(u^2-v^2\gamma)=16\gamma^2a^{q+1}$.

  We assume that $a^{q+1}\ne1$, for otherwise $a\in N$ and we are
  done. Then the product $P(X)Q(X)$ is a separable polynomial of
  degree $\le4$.

  Note that if $(x,y)$ is a point on $P(X)Y^2 = \gamma Q(X)$, then
  $P(x)\ne0$. Thus the points $(x,y)$ on the curve
  $P(X)Y^2 = \gamma Q(X)$ give different points $(x, yP(x))$ on the
  curve $Y^2=\gamma P(X)Q(X)$.

  We see that the curve $Y^2=\gamma P(X)Q(X)$ has at least $2(q-6)$
  affine points over $\FF_q$. On the other hand, by Lemma
  \ref{l:hasse}, this number is at most $q+1+2\sqrt{q}$. From this one
  quickly obtains $q\le19$.
\end{proof}

\end{document}
